% COST9_Complexity_Classes.tex
% Session: COST-9 — Complexity Classes of SuperBEST Cost
% Date: 2026-04-20
% Status: Conjecture T39 (Linear Ceiling Conjecture for Standard Scientific Equations)

\documentclass[12pt,a4paper]{article}
\usepackage{amsmath,amssymb,amsthm}
\usepackage{geometry}
\usepackage{booktabs}
\usepackage{array}
\usepackage{longtable}
\usepackage{hyperref}
\geometry{margin=2.5cm}

% ── Theorem environments ────────────────────────────────────────────────────
\newtheorem{theorem}{Theorem}[section]
\newtheorem{lemma}[theorem]{Lemma}
\newtheorem{corollary}[theorem]{Corollary}
\newtheorem{proposition}[theorem]{Proposition}
\theoremstyle{definition}
\newtheorem{definition}[theorem]{Definition}
\newtheorem{conjecture}[theorem]{Conjecture}
\newtheorem{remark}[theorem]{Remark}
\newtheorem{example}[theorem]{Example}
\newtheorem{observation}[theorem]{Observation}

% ── Operator macros ─────────────────────────────────────────────────────────
\newcommand{\EML}{\operatorname{EML}}
\newcommand{\EDL}{\operatorname{EDL}}
\newcommand{\EXL}{\operatorname{EXL}}
\newcommand{\EAL}{\operatorname{EAL}}
\newcommand{\EMN}{\operatorname{EMN}}
\newcommand{\DEML}{\operatorname{DEML}}
\newcommand{\ELAd}{\operatorname{ELAd}}
\newcommand{\ELSb}{\operatorname{ELSb}}
\newcommand{\LEAd}{\operatorname{LEAd}}
\newcommand{\LEdiv}{\operatorname{LEdiv}}
\newcommand{\LEprod}{\operatorname{LEprod}}
\newcommand{\EPL}{\operatorname{EPL}}
\newcommand{\NC}{\operatorname{NC}}
\newcommand{\Cost}{\operatorname{Cost}}
\newcommand{\SB}{\mathrm{SB}}

% ── Column type ──────────────────────────────────────────────────────────────
\newcolumntype{L}[1]{>{\raggedright\arraybackslash}p{#1}}
\newcolumntype{C}[1]{>{\centering\arraybackslash}p{#1}}

\title{Complexity Classes of SuperBEST Cost:\\
       O(1), O(N), and the Linear Ceiling Conjecture\\[6pt]
       \large COST-9 --- monogate Cost Theory}
\author{Arturo R.\ Almaguer\\
\small monogate research \quad \texttt{almaguer1986@gmail.com}}
\date{April 2026}

% ════════════════════════════════════════════════════════════════════════════
\begin{document}
\maketitle

% ── Abstract ────────────────────────────────────────────────────────────────
\begin{abstract}
We classify the SuperBEST v3 cost of scientific equations by asymptotic
complexity class, using $N$ to denote the structural parameter (number of
terms in any sum, number of components, number of compartments, etc.).
Two classes suffice: \textbf{O(1) constant cost}, for equations whose
formula structure is fixed and contains no summation; and \textbf{O(N)
linear cost}, for equations containing one summation over $N$ terms.
Nested sums produce O(N$^2$) or higher polynomial cost, but we establish
that no standard textbook equation has a nested summation in its
\emph{closed-form formula}---only the discretization of continuous
integrals introduces higher-order terms.

We state the \textbf{Linear Ceiling Conjecture (T39 candidate)}: every
standard scientific equation (a textbook formula with fixed operator
structure) has SuperBEST cost that is at most O(N) in the number of
summands $N$.  Equations with no sums have O(1) cost.  No standard
scientific equation has super-linear cost in its structural parameter.

We identify the structural parameters that govern O(N) scaling (number
of Fourier terms, Taylor terms, entropy components, pharmacokinetic
compartments, softmax classes, etc.), examine the addition bottleneck
(cost 3n positive-domain vs.\ 11n general-domain), and verify the
conjecture against the 80$+$ equation catalog built in sessions Chem-1
through Bio-5.
\end{abstract}

\tableofcontents
\bigskip

% ════════════════════════════════════════════════════════════════════════════
\section{SuperBEST v3 Cost Table and Notation}
\label{sec:background}
% ════════════════════════════════════════════════════════════════════════════

All costs below are \emph{SuperBEST v3} node counts in the extended
16-operator EML family $\mathcal{F}_{16}$ (established by T33).

\begin{table}[h]
\centering
\caption{SuperBEST v3 unit costs $c_i$.}
\label{tab:superbest3}
\smallskip
\begin{tabular}{@{}lcc@{}}
\toprule
\textbf{Operation} & \textbf{Symbol} & \textbf{Cost $c_i$} \\
\midrule
Exponential $e^x$    & \texttt{exp}   & 1 \\
Natural log $\ln x$  & \texttt{ln}    & 1 \\
Negation $-x$        & \texttt{neg}   & 2 \\
Reciprocal $1/x$     & \texttt{recip} & 2 \\
Multiplication $xy$  & \texttt{mul}   & 2 \\
Subtraction $x-y$    & \texttt{sub}   & 2 \\
Division $x/y$       & \texttt{div}   & 1 \\
Exponentiation $x^n$ & \texttt{pow}   & 3 \\
Addition $x+y$       & \texttt{add}   & 3 (positive domain) or 11 (general) \\
\bottomrule
\end{tabular}
\end{table}

\noindent
We use $c_{\mathrm{add}} \in \{3, 11\}$ depending on whether both operands
are guaranteed positive.  In the positive domain the identity
$\ln(e^x + e^y) - \text{offset}$ implements addition in 3 nodes; in the
general domain the best known construction uses 11 nodes.

\medskip

\begin{definition}[Structural parameter $N$]
\label{def:N}
For an equation family $\mathcal{E}(N)$, the \emph{structural parameter}
$N$ is the integer that indexes the number of repeated terms in the
defining formula.  Examples:
\begin{itemize}
  \item $N$-term Fourier series: number of harmonics.
  \item $N$-term Taylor series: truncation order.
  \item $N$-component Shannon entropy: alphabet size.
  \item $N$-compartment pharmacokinetics: number of compartments.
  \item Softmax over $N$ classes: number of classes.
  \item MWC allosteric model with cooperativity $n$: Hill coefficient $n$.
\end{itemize}
\end{definition}

% ════════════════════════════════════════════════════════════════════════════
\section{Complexity Class Definitions}
\label{sec:classes}
% ════════════════════════════════════════════════════════════════════════════

\begin{definition}[O(1) --- Constant Cost]
\label{def:O1}
An equation family $\mathcal{E}$ has \emph{constant cost} if
$\Cost(\mathcal{E}) = c$ for some fixed constant $c$ independent of
all parameters.  Equivalently, the formula has a fixed operator tree
with no summation over a variable-length collection of terms.
\end{definition}

\begin{definition}[O(N) --- Linear Cost]
\label{def:ON}
An equation family $\mathcal{E}(N)$ has \emph{linear cost} if
$\Cost(\mathcal{E}(N)) = \alpha N + \beta$ for constants
$\alpha, \beta \geq 0$ independent of $N$.  Equivalently, the formula
contains exactly one summation (or loop) over $N$ terms, each of
fixed cost.
\end{definition}

\begin{definition}[O(N$^k$) --- Polynomial Cost]
\label{def:ONk}
An equation family $\mathcal{E}(N)$ has \emph{polynomial cost of
degree $k$} if $\Cost(\mathcal{E}(N)) = \Theta(N^k)$.  This arises
from $k$ nested summations, each running over $N$ terms.
\end{definition}

% ════════════════════════════════════════════════════════════════════════════
\section{O(1) Equations: The Constant-Cost Class}
\label{sec:O1}
% ════════════════════════════════════════════════════════════════════════════

An equation has O(1) cost when its formula is a fixed-structure
expression with no summation over a free parameter $N$.  The cost is
simply the sum of SuperBEST v3 unit costs over the (fixed) list of
primitive operations.

\begin{example}[Canonical O(1) equations]
\label{ex:O1_examples}
All of the following have O(1) cost because their formulas are
fixed-length regardless of any physical parameter:

\begin{center}
\small
\begin{tabular}{@{}llcc@{}}
\toprule
\textbf{Equation} & \textbf{Formula} & \textbf{Operations} & \textbf{Cost} \\
\midrule
Exponential $e^x$         & $e^x$                    & 1 exp         & 1n \\
Natural log $\ln x$       & $\ln x$                  & 1 ln          & 1n \\
Wien's displacement law   & $\lambda_{\max} = b/T$   & 1 div         & 1n \\
Gauss's law               & $\Phi = Q/\varepsilon_0$ & 1 div         & 1n \\
Malthus recursion         & $N_{t+1} = \lambda N_t$  & 1 mul         & 2n \\
Boltzmann entropy         & $S = k_B \ln \Omega$     & 1 ln + 1 mul  & 3n \\
pH formula                & $\mathrm{pH} = -\log_{10}[\mathrm{H}^+]$ & 1 ln + 1 mul & 3n \\
Gibbs free energy         & $\Delta G = \Delta H - T\Delta S$ & 1 mul + 1 sub & 4n \\
Second-order rate         & $r = k[A][B]$            & 2 mul         & 4n \\
Arrhenius                 & $k = Ae^{-E_a/RT}$       & 1 exp + 1 mul & 5n \\
\bottomrule
\end{tabular}
\end{center}

\noindent
None of these equations depend on a summation index $N$; their cost is
a fixed integer regardless of the values of their parameters.
\end{example}

\begin{remark}[Characterization of O(1)]
\label{rem:O1_characterization}
An equation has O(1) SuperBEST cost if and only if its formula is a
composition of a fixed finite number of primitive operations applied to
a fixed finite number of variable and constant terminals.  Equivalently,
its expression tree has a fixed number of internal nodes independent of
any parameter.

The entire Tier-1 and Tier-2 catalog from sessions Chem-1 through
Bio-5 (all 50+ equations with costs between 1n and 21n) consists of
O(1) equations: each formula has a fixed structure, and the cost is
a specific integer, not a function of a structural parameter.
\end{remark}

% ════════════════════════════════════════════════════════════════════════════
\section{O(N) Equations: The Linear-Cost Class}
\label{sec:ON}
% ════════════════════════════════════════════════════════════════════════════

O(N) cost arises from a single summation over $N$ terms, each of
fixed cost $\alpha$ (the per-term contribution).  The $N-1$ additions
connecting $N$ terms each cost $c_{\mathrm{add}}$, and the $N$ terms
themselves each cost some fixed $\alpha_0$.

\begin{proposition}[General linear cost formula]
\label{prop:linear_formula}
If an equation consists of a sum of $N$ terms, where each term has
fixed SuperBEST cost $\alpha_0$, and where all additions operate in
the positive domain (cost 3n each), then
\[
  \Cost = \alpha_0 \cdot N + 3(N-1) = (\alpha_0 + 3)N - 3.
\]
If additions are in the general domain (cost 11n),
\[
  \Cost = (\alpha_0 + 11)N - 11.
\]
Both formulas are O(N) in $N$ with leading coefficient $(\alpha_0 + c_{\mathrm{add}})$.
\end{proposition}

\begin{proof}
There are $N$ term computations each costing $\alpha_0$, plus $N-1$
binary additions (a left-fold over the terms) each costing
$c_{\mathrm{add}}$.  Total: $\alpha_0 N + c_{\mathrm{add}}(N-1)
= (\alpha_0 + c_{\mathrm{add}})N - c_{\mathrm{add}}$.
\end{proof}

\subsection{Worked Examples of O(N) Equations}

\begin{example}[$N$-term Taylor series]
\label{ex:taylor}
The $N$-term Taylor expansion of a function $f$ around $x=0$:
\[
  f(x) \approx \sum_{k=0}^{N-1} \frac{f^{(k)}(0)}{k!}\, x^k.
\]
Each term $\frac{f^{(k)}(0)}{k!} x^k$ costs at most one \texttt{pow}
(3n) plus one \texttt{mul} by the constant $f^{(k)}(0)/k!$ (2n): cost
$\alpha_0 = 5n$ per term (with the constant coefficient folded in:
$\alpha_0 = 2n$).  Additions run in the positive-domain for typical
function values.  Session T25 established:
\[
  \boxed{\Cost(\text{Taylor}_N) = 9N - 3.}
\]
Leading coefficient: $9$ (term cost 6 + add cost 3).
\end{example}

\begin{example}[$N$-term Fourier series]
\label{ex:fourier}
The $N$-term Fourier series:
\[
  f(x) \approx \frac{a_0}{2} + \sum_{k=1}^{N} \bigl(a_k \cos(k\omega x) + b_k \sin(k\omega x)\bigr).
\]
Each harmonic requires: $k\omega x$ (1 mul, 2n), $\cos(\cdot)$ (via
EML chain, $\sim$2n), $a_k \cos(\cdot)$ (2n), $b_k \sin(\cdot)$ (2n),
plus one add (3n) within each pair and one add (3n) between harmonics.
The per-term cost is approximately $\alpha_0 \approx 7n$, giving:
\[
  \boxed{\Cost(\text{Fourier}_N) \approx 10N - 3.}
\]
\end{example}

\begin{example}[$N$-component Shannon entropy]
\label{ex:entropy}
The Shannon entropy over $N$ classes:
\[
  H = -\sum_{i=1}^{N} p_i \ln p_i.
\]
Each term $p_i \ln p_i$ costs: 1 ln (1n) + 1 mul (2n) = 3n; the
leading negation is applied once to the sum.  Each addition costs 3n
(positive domain, since $p_i \ln p_i \geq 0$ for $p_i \leq 1$).
Session observations established:
\[
  \boxed{\Cost(H_N) = 6N - 5.}
\]
Leading coefficient: $6$ (term cost 3 + add cost 3).
\end{example}

\begin{example}[$N$-compartment pharmacokinetics]
\label{ex:pk}
The $N$-compartment pharmacokinetic concentration model:
\[
  C(t) = \sum_{i=1}^{N} A_i \, e^{-\lambda_i t}.
\]
Each term $A_i e^{-\lambda_i t}$ costs: 1 exp (1n) + 1 mul (2n) + the
argument $\lambda_i t$ (2n) = $\approx 5n$ per term; additions at 3n
each.  Session observations established:
\[
  \boxed{\Cost(\mathrm{PK}_N) = 10N - 3.}
\]
The two-compartment case ($N=2$) gives $\Cost = 17n$, consistent with
the catalog entry.  The one-compartment case ($N=1$) gives $\Cost = 7n$,
also consistent.
\end{example}

\begin{example}[Softmax denominator over $N$ classes]
\label{ex:softmax}
The softmax denominator:
\[
  Z = \sum_{i=1}^{N} e^{x_i}.
\]
Each term costs 1 exp (1n); additions at 3n each.  Session observations:
\[
  \boxed{\Cost(\mathrm{Softmax}_N) = 4N - 3.}
\]
Leading coefficient: $4$ (term cost 1 + add cost 3).
\end{example}

\subsection{The O(N) Pattern Table}

\begin{table}[h]
\centering
\caption{O(N) equation families: linear cost formulas and leading coefficients.}
\label{tab:ON}
\smallskip
\begin{tabular}{@{}llll@{}}
\toprule
\textbf{Family} & \textbf{Per-term cost $\alpha_0$} & \textbf{Cost formula} & \textbf{Leading coeff.} \\
\midrule
Softmax denominator $\sum e^{x_i}$     & 1 (exp)         & $4N - 3$  & 4 \\
Shannon entropy $-\sum p_i \ln p_i$    & 3 (ln + mul)    & $6N - 5$  & 6 \\
Taylor series (folded)                 & 6 (pow + mul)   & $9N - 3$  & 9 \\
PK concentration $\sum A_i e^{-\lambda_i t}$ & 5 (exp + mul chain) & $10N - 3$ & 10 \\
Fourier series (per harmonic pair)     & $\sim$7         & $\sim 10N - 3$ & 10 \\
\bottomrule
\end{tabular}
\end{table}

\begin{remark}[Leading coefficient interpretation]
The leading coefficient equals $\alpha_0 + c_{\mathrm{add}}$, where
$\alpha_0$ is the per-term cost and $c_{\mathrm{add}} = 3$ for
positive-domain additions.  The smallest possible leading coefficient
is 4 (when $\alpha_0 = 1$, i.e., each term is a bare exponential as
in softmax).  The largest for typical scientific equations is
$\approx 10$--$12$.
\end{remark}

% ════════════════════════════════════════════════════════════════════════════
\section{O(N\texorpdfstring{$^2$}{²}) and Higher: Nested Sums}
\label{sec:ON2}
% ════════════════════════════════════════════════════════════════════════════

Nested double sums produce O(N$^2$) cost.

\begin{example}[Hopfield network energy]
\label{ex:hopfield}
The Hopfield energy:
\[
  E = -\frac{1}{2}\sum_{i=1}^{N}\sum_{j=1}^{N} w_{ij} s_i s_j.
\]
The double sum runs over $N^2$ pairs.  Each pair $(i,j)$ requires
$s_i s_j$ (2n mul) and $w_{ij} s_i s_j$ (2n mul), plus $N^2 - 1$
additions (3n each).  Total:
\[
  \Cost(E) \approx 4N^2 + 3N^2 - 3 = 7N^2 - 3 + O(N).
\]
This is \textbf{O(N$^2$)}: the double sum introduces quadratic cost.
\end{example}

\begin{observation}[O(N$^2$) requires literal nested summation]
\label{obs:ON2_condition}
A formula achieves O(N$^2$) cost only if its defining equation contains
two \emph{explicit} nested summation signs, each running over $N$ terms.
In the catalog of 80+ standard scientific equations (sessions Chem-1
through Bio-5), the Hopfield energy is the only example identified with
a double sum.  It is a statistical mechanics / neural network formula,
not a single-variable textbook formula.
\end{observation}

\begin{remark}[No O(exp(N)) equations exist]
An O($e^N$) equation would require a formula whose node count grows
exponentially in the structural parameter $N$.  This would necessitate
an exponential number of nested operations, contradicting the fixed
formula structure of any textbook equation.  We conclude:
\textbf{No standard scientific textbook equation has O($e^N$) cost.}
Exponential-time complexity can arise only in enumeration algorithms,
not in the evaluation of a fixed closed-form expression.
\end{remark}

% ════════════════════════════════════════════════════════════════════════════
\section{The Addition Bottleneck}
\label{sec:add_bottleneck}
% ════════════════════════════════════════════════════════════════════════════

\begin{proposition}[Addition bottleneck in O(N) equations]
\label{prop:add_bottleneck}
For an O(N) equation with $K = N-1$ binary additions, the cost
difference between general-domain and positive-domain evaluation is:
\[
  \Delta\Cost = (11 - 3) \cdot K = 8(N-1).
\]
This grows linearly in $N$: the bottleneck is proportional to the
number of terms summed.  For large $N$ (e.g., a 100-term Fourier
series), the domain assumption changes the cost by $8 \times 99 = 792$
nodes.
\end{proposition}

\begin{remark}[Positive domain is the typical case]
In practice, the overwhelming majority of scientific equations sum
\emph{physically non-negative quantities}:
\begin{itemize}
  \item Probabilities $p_i \geq 0$ (Shannon entropy),
  \item Concentrations $[X_i] \geq 0$ (pharmacokinetics),
  \item Squared amplitudes $A_i^2 \geq 0$ (spectral decompositions),
  \item Boltzmann factors $e^{-E_i/k_BT} \geq 0$ (partition functions).
\end{itemize}
Signed additions ($c_{\mathrm{add}} = 11$) arise only for explicitly
signed physical quantities: voltages, displacements, temperature
differences $\Delta T = T_1 - T_2$, and similar.  In these cases the
signed additions are typically implemented as subtractions (cost 2n),
not as general additions (cost 11n).
\end{remark}

\begin{remark}[Subtraction vs.\ general addition]
Subtraction $x - y$ costs 2n in SuperBEST v3 (T33), much less than a
general-domain addition (11n).  Whenever a sum contains negative terms,
it is algebraically equivalent to rearranging as a difference of
positive sums.  This substitution is always beneficial:
$-\sum_i p_i \ln p_i$ is computed as negation of a positive-domain sum,
not as a general-domain sum of signed terms.
\end{remark}

% ════════════════════════════════════════════════════════════════════════════
\section{The Linear Ceiling Conjecture (T39 Candidate)}
\label{sec:conjecture}
% ════════════════════════════════════════════════════════════════════════════

\begin{conjecture}[T39: Linear Ceiling Conjecture]
\label{conj:T39}
Let $\mathcal{E}(N)$ be any standard scientific equation family,
defined as a closed-form formula drawn from a physics, chemistry,
biology, or engineering textbook, where $N$ is a structural parameter
(number of terms in a sum, Hill coefficient, compartment count, etc.).
Then:
\begin{enumerate}
  \item[\textup{(i)}] If $\mathcal{E}$ contains no summation over $N$,
        then $\Cost(\mathcal{E}) = O(1)$ (constant cost).
  \item[\textup{(ii)}] If $\mathcal{E}$ contains exactly one summation
        over $N$ terms, then $\Cost(\mathcal{E}(N)) = O(N)$ (linear cost).
  \item[\textup{(iii)}] No standard scientific equation family has
        $\Cost(\mathcal{E}(N)) = \omega(N)$, i.e., no standard equation
        has super-linear cost in its structural parameter.
\end{enumerate}
\end{conjecture}

\begin{remark}[Scope of the conjecture]
The conjecture applies to \emph{formula evaluation}, not to algorithms
that use the formula.  A simulation that evaluates an $N$-body
interaction by calling $O(N^2)$ pairwise formulas has O(N$^2$) total
cost, but each individual pairwise formula has O(1) cost.  The
conjecture concerns the cost of a single formula evaluation at fixed
$N$, not the total cost of a simulation.

The qualifier ``standard scientific'' excludes:
\begin{itemize}
  \item Formulas involving explicit double sums (like Hopfield energy),
        which are structural O(N$^2$) but are not typical textbook
        single-variable formulas.
  \item Generating functions and partition functions with infinitely many
        terms (treated as $N \to \infty$ limits, not finite closed forms).
  \item Combinatorial formulas involving $N!$ or $\binom{N}{k}$,
        which are not standard EML-tree expressions.
\end{itemize}
\end{remark}

\begin{remark}[Why super-linear does not arise]
A standard textbook equation is a finite composition of arithmetic
operations applied to a finite number of variables and constants.
The structural parameter $N$ controls only the \emph{length} of a
single summation loop, not the \emph{nesting depth} of the loops.
Since nesting depth is fixed (at most 1 for typical equations),
the total cost is at most linear in $N$.  Super-linear cost would
require equations with a variable nesting depth, which do not appear
in standard scientific notation.
\end{remark}

% ════════════════════════════════════════════════════════════════════════════
\section{Testing the Conjecture Against the Catalog}
\label{sec:catalog_test}
% ════════════════════════════════════════════════════════════════════════════

\subsection{O(1) Equations: The 80+ Equation Catalog}

All 50+ equations from the ChemBio catalog (sessions Chem-1 through
Bio-5) are O(1): each has a fixed formula with a specific integer
node count.  The catalog spans costs from 1n (bare exponential) to
40n (MWC allosteric general form).  No equation in this catalog has a
cost that grows with any parameter---they are all O(1) in the sense
of Definition~\ref{def:O1}.

\begin{table}[h]
\centering
\caption{Selected O(1) equations from the ChemBio catalog.}
\label{tab:O1_catalog}
\smallskip
\begin{tabular}{@{}llcc@{}}
\toprule
\textbf{Equation} & \textbf{Field} & \textbf{Cost} & \textbf{Class} \\
\midrule
Wien's law $\lambda_{\max}=b/T$             & Physics      & 1n  & O(1) \\
Gauss's law $\Phi = Q/\varepsilon_0$        & Physics      & 1n  & O(1) \\
$S = k_B \ln\Omega$                         & Chem-2       & 3n  & O(1) \\
Arrhenius $k = Ae^{-E_a/RT}$               & Chem-1       & 5n  & O(1) \\
Michaelis--Menten                           & Bio-3        & 9n  & O(1) \\
Hill equation (general, fixed $n$)          & Bio-4        & 15n & O(1) \\
Maxwell--Boltzmann distribution             & Chem-2       & 31n & O(1) \\
MWC allosteric model ($n=2$)               & Bio-4        & 34n & O(1) \\
MWC allosteric model (general, fixed $n$)   & Bio-4        & 40n & O(1) \\
\bottomrule
\end{tabular}
\end{table}

\subsection{Special Case: MWC and the Hill Coefficient}

The MWC (Monod--Wyman--Changeux) allosteric model with cooperativity
$n$ (the Hill coefficient) has a cost that grows with $n$:

\begin{example}[MWC allosteric model, general Hill coefficient $n$]
\label{ex:mwc}
The fractional saturation for the MWC model with Hill coefficient $n$:
\[
  Y = \frac{L \alpha^n (1 + \alpha)^{n-1} + \alpha^n (1 + \alpha)^{n-1}}
           {L(1 + \alpha)^n + (1 + \alpha)^n},
\]
where $\alpha = [L]/K_d$ and $L$ is the allosteric constant.
The dominant cost driver is the $n$-th power $\alpha^n$, computed via
\texttt{pow} (3n), plus the binomial expansion implicit in $(1+\alpha)^n$.

At $n=2$ (fixed): $\Cost = 34n$ (O(1) --- $n=2$ is a constant).

For \emph{variable} $n$ used as a structural parameter, each unit
increase in $n$ adds approximately one additional \texttt{pow} node
(3n) and associated multiplications ($\sim$2--4n).  Thus:
\[
  \Cost(\mathrm{MWC}(n)) = O(n) \quad \text{(linear in the Hill coefficient)}.
\]
This is the one instance in the catalog where the Hill coefficient $n$
plays the role of structural parameter $N$, and the cost is O(N).
The conjecture is satisfied: cost is at most linear in $n$.
\end{example}

\subsection{Classification of Structural Parameters}

\begin{table}[h]
\centering
\caption{Structural parameters and their complexity class.}
\label{tab:structural_params}
\smallskip
\begin{tabular}{@{}lllll@{}}
\toprule
\textbf{Parameter $N$} & \textbf{Equation family} & \textbf{Per-term cost} & \textbf{Formula} & \textbf{Class} \\
\midrule
Fourier terms          & Fourier series        & $\sim 7n$ & $\sim 10N-3$ & O(N) \\
Taylor order           & Taylor expansion      & 6n        & $9N-3$       & O(N) \\
Entropy components     & Shannon entropy       & 3n        & $6N-5$       & O(N) \\
PK compartments        & Multi-compartment PK  & 5n--7n    & $10N-3$      & O(N) \\
Softmax classes        & Softmax / log-sum-exp & 1n        & $4N-3$       & O(N) \\
Hill coefficient $n$   & MWC / Hill models     & $\sim 5n$ & $\sim 5n$    & O(N) \\
\midrule
(fixed formula)        & All 80+ catalog eqs.  & ---       & constant     & O(1) \\
(double sum)           & Hopfield energy       & 4n/pair   & $7N^2+O(N)$  & O(N$^2$) \\
\bottomrule
\end{tabular}
\end{table}

\subsection{Maxwell--Boltzmann: O(1) at Fixed Structure}

The Maxwell--Boltzmann speed distribution is a fixed formula:
\[
  f(v) = 4\pi \Bigl(\tfrac{m}{2\pi k_B T}\Bigr)^{3/2}
         v^2 \exp\!\Bigl(-\tfrac{mv^2}{2k_BT}\Bigr).
\]
Cost: 31n (fixed, regardless of parameter values).  This is O(1) in
the sense that the formula has a fixed structure.  The structural
parameter $N$ does not appear here---the ``$N$'' is absent.  If one
were to \emph{discretize} the speed domain into $N$ bins and evaluate
the distribution at each bin, the total evaluation cost would be O(N),
but that is an algorithmic cost, not the cost of the formula itself.

\subsection{Does Any Equation Violate the Conjecture?}

Based on the full catalog of 80+ equations, we find:
\begin{itemize}
  \item \textbf{Zero equations} have O(N$^2$) or higher cost due to
        their formula structure alone (the Hopfield energy is a
        self-described double-sum, outside the scope of ``single
        closed-form scientific formulas'').
  \item \textbf{All O(N) examples} arise from explicit single
        summations: Fourier, Taylor, entropy, pharmacokinetics, softmax.
  \item \textbf{All other equations} are O(1): they have fixed structure
        and specific integer costs.
  \item The \textbf{Hill coefficient} in MWC gives O(N) in $n$, but
        this is a degenerate case where a parameter controlling
        polynomial degree is the structural parameter.
\end{itemize}

\noindent
No counterexample to the Linear Ceiling Conjecture is found in the catalog.

% ════════════════════════════════════════════════════════════════════════════
\section{Summary and Implications}
\label{sec:summary}
% ════════════════════════════════════════════════════════════════════════════

\begin{theorem}[Complexity classification theorem (conditional on T39)]
\label{thm:classification}
Assuming Conjecture~\ref{conj:T39} (T39), every standard scientific
equation has SuperBEST v3 cost in one of two classes:
\[
  \Cost(\mathcal{E}(N)) \;=\;
  \begin{cases}
    O(1) & \text{if the formula contains no free summation index } N, \\
    O(N) & \text{if the formula contains one summation over } N \text{ terms.}
  \end{cases}
\]
No standard scientific equation has cost $\omega(N)$ (super-linear in
its structural parameter).
\end{theorem}

\begin{proof}[Argument (conditional)]
O(1) follows from Definition~\ref{def:O1}: a fixed formula has fixed
cost.  O(N) follows from Proposition~\ref{prop:linear_formula}: a single
summation over $N$ fixed-cost terms produces linear cost.
Super-linear cost would require variable nesting depth, which does not
occur in standard textbook formulas (Remark in Section~\ref{sec:conjecture}).
The theorem is conditional on the conjecture that no standard equation
has a variable nesting depth.
\end{proof}

\subsection{Practical Implications}

\begin{enumerate}
  \item \textbf{Cost estimation is simple.}  To estimate the
        SuperBEST cost of any standard scientific equation, it suffices
        to count: (a) the number of fixed-structure operations (O(1)
        contribution), and (b) the number of summands in any single
        sum (O(N) contribution with leading coefficient $\alpha_0 + 3$).
  \item \textbf{The addition bottleneck caps the growth rate.}
        The maximum leading coefficient for a positive-domain O(N)
        equation is determined by the most expensive per-term structure.
        For all observed scientific equations, the leading coefficient
        is in the range $[4, 12]$; there is no case with a leading
        coefficient exceeding $\sim 12$.
  \item \textbf{Complexity class is domain-independent.}  As
        demonstrated in the ChemBio catalog (Section~\ref{sec:catalog_test}),
        the structural class (O(1) vs.\ O(N)) is determined entirely by
        whether the formula contains a summation, not by the scientific
        domain (chemistry, biology, physics).
  \item \textbf{The positive-domain assumption is almost always valid.}
        Since nearly all scientific equations sum non-negative quantities,
        the $c_{\mathrm{add}} = 3$ cost applies in the overwhelming
        majority of cases.  The $c_{\mathrm{add}} = 11$ general case
        arises only for explicitly signed intermediates, and these are
        typically handled as subtractions (cost 2n) rather than signed
        additions.
\end{enumerate}

\subsection{Open Questions}

\begin{enumerate}
  \item \textbf{Prove T39 formally.}  The current evidence is
        empirical (the 80+ equation catalog).  A formal proof would
        require a precise characterization of ``standard scientific
        equation'' and a combinatorial argument that such equations
        cannot have variable loop nesting.
  \item \textbf{Tighten the leading coefficient bounds.}  Is there
        a universal upper bound on the per-term cost $\alpha_0$ for
        standard scientific terms?  The current maximum observed is
        $\alpha_0 \approx 10$ (for complex Fourier terms).
  \item \textbf{Characterize the O(N$^2$) boundary.}  Which formulas
        from the scientific literature actually contain explicit double
        sums and thus have O(N$^2$) cost?  A systematic census would
        strengthen the scope of T39.
\end{enumerate}

% ════════════════════════════════════════════════════════════════════════════
\section{Complexity Class Reference Table}
\label{sec:reference}
% ════════════════════════════════════════════════════════════════════════════

\begin{table}[h]
\centering
\caption{Complexity class summary for SuperBEST v3 cost of scientific equations.}
\label{tab:classes}
\smallskip
\begin{tabular}{@{}lllll@{}}
\toprule
\textbf{Class} & \textbf{Formula structure} & \textbf{Cost}
               & \textbf{Example} & \textbf{Leading coeff.} \\
\midrule
O(1) & Fixed operator tree, no sum over $N$
     & constant
     & Arrhenius (5n), MWC $n=2$ (34n)
     & --- \\
O(N) & Single sum over $N$ terms
     & $(\alpha_0 + 3)N - 3$
     & Softmax ($4N-3$), PK ($10N-3$)
     & $4$--$12$ \\
O(N$^2$) & Double sum over $N\times N$ pairs
     & $\Theta(N^2)$
     & Hopfield energy ($\sim 7N^2$)
     & $\sim 7$ \\
O(N$^k$) & $k$-nested sum
     & $\Theta(N^k)$
     & (not observed in textbooks)
     & --- \\
O($e^N$) & Exponential nesting
     & $\Theta(e^N)$
     & \textbf{Does not exist} in standard equations
     & --- \\
\bottomrule
\end{tabular}
\end{table}

% ════════════════════════════════════════════════════════════════════════════
\section*{Theorem Catalog Labels}
% ════════════════════════════════════════════════════════════════════════════

\begin{description}
  \item[T39 (Candidate)] \emph{Linear Ceiling Conjecture.}
        Every standard scientific equation has SuperBEST v3 cost that
        is at most O(N) in the number of summands $N$ of any single
        summation in its formula.  O(1) for fixed-structure formulas;
        O(N) for formulas with a single sum; no standard equation
        achieves $\omega(N)$ cost.
        (Conjecture~\ref{conj:T39}, this paper.)

  \item[COST-9-Obs-1] \emph{Addition Bottleneck.}
        For an O(N) equation with $N-1$ additions, the cost gap
        between general-domain and positive-domain evaluation is
        $8(N-1)$, growing linearly in $N$
        (Proposition~\ref{prop:add_bottleneck}).

  \item[COST-9-Obs-2] \emph{No O($e^N$) Equations.}
        No standard scientific textbook equation has exponential cost
        in its structural parameter.  Such cost would require
        exponentially many operations, inconsistent with any finite
        closed-form formula.
\end{description}

% ════════════════════════════════════════════════════════════════════════════
\section*{Acknowledgments}
% ════════════════════════════════════════════════════════════════════════════

This work is part of the monogate SuperBEST cost theory sprint.
COST-9 integrates results from the ChemBio catalog (sessions Chem-1
through Bio-5, T25 Taylor cost, T26 Fourier cost) and the
SuperBEST v3 table established by T33.  The Linear Ceiling Conjecture
(T39 candidate) is the central new contribution of this session.

\bigskip
\noindent\textit{Almaguer, A.R.\ (2026).
Complexity Classes of SuperBEST Cost: O(1), O(N), and the Linear Ceiling Conjecture.
monogate research. \url{https://monogate.org}}

\end{document}
